\chapter{Selected Hints and Short Answers}

This appendix gathers the style of answers used throughout the book: enough to check the main idea, not enough to replace doing the work.

\section{Differentiation under an integral}

When a source \(J\) appears in an exponential, differentiating with respect to \(J(x)\) brings down the field at \(x\):
\[
  \frac{\delta}{\delta J(x)}
  e^{\ii\int J\phi}
  =
  \ii\phi(x)e^{\ii\int J\phi}.
\]
This is why generating functionals generate correlation functions.

\section{Boundary terms}

Most variational derivations in the text drop boundary terms.  The usual reasons are:
\[
  \delta\phi=0\quad\text{on the boundary},
  \qquad
  \phi\to0\quad\text{fast enough at infinity},
\]
or periodic boundary conditions.  If none of these holds, the boundary term is physical and must be kept.

\section{Checking signs}

When signs in propagators become confusing, return to the defining inverse equation.  With the conventions of this book,
\[
  (\Boxop+m^2)\Delta_F(x)=-\ii\delta^{(4)}(x).
\]
For fermions,
\[
  (\ii\gamma^\mu\p_\mu-m)S_F(x)=\ii\delta^{(4)}(x).
\]
These equations determine the signs in momentum space.

